\begin{frame}{PSUM = TLE}
    \metroset{block=fill}
    
    \begin{center}
        \tikz {
        \node (1) at (0,0) {$v_1$};
        \node (2) at (1,0) {$v_2$};
        \node (3) at (2,0) {$v_3$};
        \node (4) at (3,0) {$v_4$};

        % \node (p1) at (0,1) {$p_1$};
        \node (p2) at (0.5,1) {$p_2$};
        % \node (p3) at (2,2) {$p_3$};
        \node (p4) at (2.5,1) {$p_4$};

        \node (r) at (1.5,2) {$r$};

        \graph { 
            (p4) -> {(4), (3)};
            % (p3) -> {(3), (p2)};
            (p2) -> {(2), (1)};
            % (p1) -> (1);
            (r) -> {(p2), (p4)};
        };
        }
    \end{center}

    Diminuímos a quantidade máxima de dependências de $4$ para $2$. Notemos que a estrutura que fizemos foi a de uma árvore binária.

    Problema: como fazer a soma $v_2+v_3+v_4$? Podemos soma $v_2$ com $p_4$.

    Será que podemos fazer uma estrutura similar para qualquer tamanho de vetor?

    % \begin{block}{Implementação da busca binária}
    %     \footnotesize
    %     \inputminted{cpp}{codigos/busca.binaria.cpp}    
    % \end{block}

\end{frame}
